\documentclass[11pt,a4paper,titlepage]{article}
%=============================================
%Preámbulo

%Idioma
\usepackage[utf8x]{inputenc}
\usepackage[spanish]{babel}

%Imgágenes y flotantes
\usepackage{graphicx}
\usepackage{float}

%Márgenes
\usepackage[a4paper,includeheadfoot,margin=2.54cm]{geometry}
%\usepackage[a4paper,margin=2.54cm]{geometry}

%Encabezado y pie de página
\usepackage{fancyhdr}

\pagestyle{fancy}
\fancyhf{}
\fancyhead[R]{\includegraphics[width=3cm]{imagenes/Logo_FIUBA.png}}
\fancyhead[L]{\includegraphics[width=3cm]{imagenes/Logo_Labi_color.png}}
\setlength{\headheight}{46pt}
\fancyfoot[R]{\thepage}
\fancyfoot[C]{\leftmark}
\renewcommand{\headrulewidth}{2pt}
\renewcommand{\footrulewidth}{1pt}

\author{LABi}
\title{Curso de KiCad - v5.1.5}
\date{1C 2020}
%=============================================
%Inicio del documento

\begin{document}
\maketitle

\tableofcontents
\newpage

\section{¿Qué es KiCad?}

Es una colección de programas open source, que en conjunto permiten diseñar un PCB, pasando por diferentes etapas.

\begin{itemize}
	\item KiCad - administrador de proyectos
	\item Eeschema - editor de esquemáticos
	\item Cvpcb - asociación de huellas (\textit{footprints}) de los componentes usados en el esquemático
	\item Pcbnew - entorno de diseño de los circuitos impresos (PCB)
	\item Gerbview - visualizador de archivos Gerber, usados para la fabricación de la placa
\end{itemize}




\end{document}
